\begin{figure}[H]
\centering
\begin{tikzpicture}[font=\small,>=Latex,
  gpu/.style={draw=gray!65,fill=gray!5,text width=2.05cm,minimum height=.85cm,align=center}]
\node[draw=blue!65!black,fill=blue!5,rounded corners,align=center,
      text width=3.6cm,minimum height=1.1cm] (p) at (0,0)
{Um processo\\Uma região \texttt{target}};
\node[gpu,draw=orange!75!black,fill=orange!18] (a) at (5.7,1.5) {Dispositivo padrão};
\node[gpu] (b) at (9.2,1.5) {Outro dispositivo};
\node[gpu] (c) at (5.7,-.3) {Outro dispositivo};
\node[gpu] (d) at (9.2,-.3) {Outro dispositivo};
\draw[->,thick] (p.east) -- ++(.6,0) |- (a.west);
\node[draw=gray!60,dashed,fit=(a)(b)(c)(d),inner sep=8pt,
      label={[font=\footnotesize]below:Exemplo: quatro GPUs visíveis na alocação}] {};
\end{tikzpicture}
\caption{Oito GPUs no nó e quatro GPUs solicitadas ao Slurm não significam uma soma distribuída por quatro dispositivos. O programa possui uma região target, sem particionamento explícito do vetor entre GPUs.}
\label{fig:dispositivos}
\end{figure}
